\chapter{Introduction}
\label{chapter:introduction}

\section{Motivation}
\label{section:motivation}

Robotics exist to serve people and help automate repetitive tasks, but most importantly, to automate tasks that are hazardous for humans. Such applications can be found in highly radioactive environments. On Earth, these are usually artificial radiation environments, such as nuclear power plants, nuclear research centers, and particle accelerators. Therefore, applications such as nuclear decommissioning, nuclear waste management, and interventions in radioactive areas require robust solutions that are tolerant to radiation~\cite{DiCastro2018CERNTAURO, Tsiligiannis2020Evaluation}.

Radiation adds a significant layer of complexity to modern electronics, as it can introduce errors, performance degradation, and permanent failures, making many commercial systems unsuitable for harsh environments. Similar challenges are encountered in other fields, such as space missions, where radiation-hardened components, often referred to as space- or military-grade devices, have been developed to operate under radiation, vacuum, and thermal stress. However, these solutions are not always viable when electronics must be deployed in large quantities, due to cost, availability, and procurement constraints. For this reason, COTS components are widely used at CERN and in other organizations, as they reduce production cost and simplify prototyping and procurement. Their use, however, is not arbitrary: depending on the radiation levels expected at each location in the accelerator complex, COTS-based systems must undergo thorough testing and qualification, typically through component-level radiation tests at dedicated facilities. This approach is supported by CERN’s Radiation Hardness Assurance (RHA) activities and by the Radiation to Electronics (R2E) project, which studies radiation effects on electronics and develops strategies for their reliable operation in accelerator environments~\cite{Tsiligiannis2020Evaluation, Ferraro2020COTS, GarciaAlia2018LHC, ratopus_jacow, Coronetti2021Radiation}.

With the vision of the next accelerator, the \textbf{F}uture \textbf{C}ircular \textbf{C}ollider (\textbf{FCC}), the role of robotics becomes even more important. At CERN, robotic systems are used to perform remote inspection, measurement, maintenance, and intervention tasks in areas where direct human access is difficult, time-consuming, or unsafe. In the FCC, this need is amplified by the scale of the machine, with a proposed underground tunnel of around \SI{100}{\kilo\metre}, where long distances, radiation exposure, and availability constraints make manual interventions increasingly challenging. For this reason, a dedicated robotic system has been proposed for the FCC, aiming to automate maintenance, inspection, and emergency-handling tasks while reducing personnel exposure and supporting reliable accelerator operation~\cite{Gamper2024Robotic}.

The work is part of a cooperation between the \textbf{E}lectronics \textbf{P}roduction \& \textbf{R}adiation tolerance (\textbf{EPR}) and \textbf{M}echatronics, \textbf{R}obotics \& \textbf{O}peration (\textbf{MRO}) sections of the \textbf{C}ontrols, \textbf{E}lectronics \& \textbf{M}echatronics (\textbf{CEM}) group.

\section{Thesis Objectives Outline}
\label{section:thesis-objectives}

Considering the aforementioned, the following question is formulated:

\vspace{1em}

\begin{center}
    \begin{minipage}{0.65\textwidth}
        \centering
        \itshape
        \normalsize
        ``Can we develop radiation-tolerant robotic arms for use in CERN accelerators, even during beam time?''
    \end{minipage}
\end{center}

\vspace{1em}

The main objective of this thesis is to lay the technical foundation for answering this question. This is done by starting from scratch and delivering a prototype servo mechanism that can act as a proof of concept for future developments.

The points of interest are, of course, to gather the requirements of a motor driver system, decide on an appropriate algorithm to drive the motor, and provide real hardware and software implementations. Furthermore, a very important objective is the development of a reproducible test setup that can be installed in different facilities and allow for further radiation tests.

More specifically, this thesis aims to address the following objectives:

\textbf{Objective 1.} Design a prototype servo system based on discrete, radiation-characterized components.

\textbf{Objective 2.} Implement an advanced motor control algorithm, while considering the need for portability across platforms, scalability, and radiation-induced error mitigation techniques.

\textbf{Objective 3.} Develop a complete motor test bench that can be reproduced and used in different facilities to evaluate hardware and control degradation under irradiation.

\textbf{Objective 4.} Identify the different error modes, analyze how they affect the control behavior, and evaluate whether safety standards are violated.

\textbf{Objective 5.} Propose ideas for future work.

An initial target of \SI{500}{\gray} was set.\label{target:500gy-goal}

The final deliverables of this work include:
\begin{enumerate}
    \item A prototype motor driver board that was proved radiation tolerant up to 650 Gy ($^{60}Co$ source)
    \item A fully-custom Field-Oriented Control (FOC) Cortex-M0+ firmware developed in C++ and including the control algorithm execution and state machines, hardware abstractions, radiation-induced fault mitigation attempts, and utilities for diagnostics and maximum during-test observability
    \item A fully reproducible motor test bench, including the mechanical CAD files for integration of a hollow-shaft BLDC motor, magnet, and golden reference external LVDT resolver at the irradiation room side, and the hardware and software at control room side
    \item Analysis of the results of the irradiation tests that were performed at CHARM and CC60 facilities at CERN
\end{enumerate}

\textbf{\textit{Note:}} the radiation tolerance of the motor itself — including insulation, copper, permanent magnets, bearings, and mechanical materials — is outside the scope of this thesis. The motor is treated as a radiation-tolerant load; the driving electronics are under investigation.

\section{Thesis Structure Outline}

The thesis is structured as follows:

\textbf{Chapter 2} introduces the fundamentals of radiation effects on electronics, covering their impact on semiconductor devices, and describes the facilities used for radiation testing.

\textbf{Chapter 3} reviews the state of the art of CERN's robotic and positioning systems, discusses the motor and control-algorithm selection, and defines the corresponding hardware requirements.

\textbf{Chapter 4} describes the design and implementation of the motor control system, covering both hardware and software, and distinguishes between the different design iterations.

\textbf{Chapter 5} details the mechanical design of the test bench, the test preparation, and the proposed methodology for analyzing the radiation-test results.

\textbf{Chapter 6} reports the radiation-test results and discusses the iterative process followed to improve the radiation tolerance of the system.

\textbf{Chapter 7} summarizes the lessons learned and proposes future steps towards further improving the system performance.